\section{Modality-Specific Applications}
\label{sec:applications}

Generative model architectures require adaptation to the distinct characteristics of different biomedical data modalities. Medical imaging, electronic health records, and multi-omics data each present unique challenges in terms of dimensionality, structure, temporal dependencies, and domain constraints. This section examines how researchers have tailored generative approaches to specific data types, with particular attention to the interplay between modality characteristics and model design choices.

\subsection{Medical Imaging}

Medical imaging applications have dominated the deployment of generative models in rare disease research, accounting for 70.3\% of reviewed studies \cite{finetti2025data}. This concentration reflects both the availability of established deep learning architectures for image data and the acute need for expanded training sets in radiology and pathology where acquisition costs and patient scarcity severely limit dataset sizes.

Frid-Adar et al. applied Deep Convolutional GANs to computed tomography (CT) imaging for liver lesion classification \cite{fridadar2018gan}. Starting from 182 lesion samples across three diagnostic categories (cysts, metastases, hemangiomas), they generated synthetic 64×64 pixel regions of interest that captured characteristic texture patterns and morphological features. The generator architecture employed fractionally-strided convolutions to progressively upsample from latent noise vectors to full-resolution images, with batch normalization stabilizing training dynamics. Classical augmentation techniques (rotations, translations, scaling, flipping) expanded each original lesion to approximately 30,000 transformed samples before GAN synthesis added further diversity. The combined augmentation strategy enabled convolutional neural network training that achieved 85.7\% sensitivity and 92.4\% specificity, representing 7.1 and 4.0 percentage point improvements over classical augmentation alone.

Expert radiologist validation confirmed the diagnostic utility of synthetic lesions, though systematic evaluation protocols remain underdeveloped. Visual inspection verified that generated samples exhibited realistic tissue density values, appropriate spatial relationships between anatomical structures, and diagnostically relevant features. However, this validation approach scales poorly to large synthetic datasets and introduces subjectivity that complicates reproducibility assessment across institutions.

The application of generative models to medical imaging extends beyond CT to magnetic resonance imaging (MRI), retinal fundus photography, and histopathology. Mendes et al. documented that combining synthetic and real brain MRI data achieved 85.9\% classification accuracy with Dice score improvements of 3-15\% for segmentation tasks \cite{mendes2025synthetic}. For age-related macular degeneration, ResNet-18 architectures trained on GAN-augmented retinal images reached 85\% diagnostic accuracy, outperforming human expert consensus at 75-80\% \cite{mendes2025synthetic}. These results demonstrate that synthetic augmentation can surpass simple dataset expansion to enable superhuman performance in specific diagnostic contexts.

The choice between GAN variants depends on the specific imaging task. Conditional GANs enable generation of specific lesion types or disease stages by incorporating class labels as additional inputs. CycleGANs facilitate unpaired image-to-image translation, useful for modality conversion or artifact removal when paired training examples are unavailable. StyleGANs provide fine-grained control over generated image attributes through manipulation of latent style vectors at different resolution scales. The continued proliferation of GAN architectures specialized for different imaging modalities reflects the domain-specific nature of optimal generative model design.

\subsection{Electronic Health Records}

Electronic health record (EHR) data represented 19.5\% of reviewed applications \cite{finetti2025data}, encompassing both cross-sectional aggregated features and longitudinal event sequences. The discrete, high-dimensional, and sparse nature of EHR data necessitates distinct architectural adaptations compared to continuous-valued imaging data.

Choi et al. developed medGAN specifically for synthesizing multi-label discrete patient records \cite{choi2017generating}. The architecture combines an autoencoder for dimensionality reduction with a GAN for generating samples in the learned latent space. The autoencoder compresses binary and count variables representing diagnosis and medication codes into continuous latent vectors, enabling smoother GAN training than operating directly on discrete features. The generator incorporates shortcut connections that propagate gradients more effectively during backpropagation, while the discriminator employs minibatch averaging to combat mode collapse without hyperparameter tuning.

Testing on three EHR datasets ranging from 30,738 to 258,559 patients demonstrated that synthetic records preserved marginal distributions and co-occurrence patterns of medical codes. Disease prediction models trained on synthetic data achieved 97.77\% accuracy, with sensitivity of 95.21\% and specificity of 99.25\%. Privacy analysis through membership inference attacks and attribute disclosure metrics confirmed that synthetic records did not memorize individual training samples, enabling data sharing between institutions without residual re-identification risk.

The temporal dimension of EHR data introduces additional complexity beyond cross-sectional aggregation. Disease progression follows characteristic trajectories that synthetic data must respect to maintain clinical validity. Zhong et al. addressed this challenge through MedDiffusion, adapting diffusion probabilistic models to sequential medical records \cite{zhong2024meddiffusion}. The architecture embeds diagnosis codes and visit timing information into continuous representations processed by bidirectional LSTMs that capture longitudinal dependencies. A step-wise attention mechanism conditions synthetic visit generation on patient history, ensuring temporal consistency in disease progression patterns.

MedDiffusion outperformed GAN baselines across four health risk prediction tasks, achieving 77.88\% PR-AUC for kidney disease prediction compared to 74.92\% for the best GAN alternative. The performance gain derives from diffusion models' superior handling of discrete sequential data and their avoidance of the mode collapse that frequently afflicts GAN training on structured medical records. The conditioning mechanism proved critical---ablation studies removing the step-wise attention or the diffusion process substantially degraded performance, confirming that both components contribute essential capabilities.

The expansion of EHR datasets through synthetic generation addresses not only quantity limitations but also representation biases. Real-world EHR systems overrepresent frequently-encountered conditions and patient demographics while underrepresenting rare diseases and minority populations. Synthetic generation with appropriate conditioning can rebalance these distributions, though careful validation remains necessary to ensure that generated samples reflect realistic disease prevalence rather than artifacts of the generation process.

\subsection{Rare Genetic Diseases and Knowledge-Guided Synthesis}

Multi-omics data and rare genetic diseases present extreme scarcity challenges where even dozens of patient samples may be unavailable. These applications constitute only 5.1\% of reviewed studies \cite{finetti2025data}, yet represent critical domains where generative approaches could provide disproportionate impact. The ultra-low sample regime necessitates integration of domain knowledge beyond purely data-driven learning.

Alsentzer et al. demonstrated an alternative paradigm through SHEPHERD, a few-shot learning approach for phenotype-driven diagnosis that relies primarily on synthetic patient data \cite{alsentzer2025few}. Rather than augmenting small real datasets, SHEPHERD trains on 42,624 simulated patients representing 2,132 unique Mendelian disorders. The synthesis process grounds generation in biomedical knowledge graphs encoding relationships between phenotypes, genes, diseases, and biological pathways. Each synthetic patient comprises phenotypic features sampled from disease-specific distributions in the Orphanet rare disease database, with additional corruption and noise injection to simulate diagnostic uncertainty and incomplete clinical information.

The SHEPHERD architecture employs graph neural networks that operate directly on patient phenotype subgraphs extracted from a heterogeneous biomedical knowledge graph containing 105,220 nodes and over 1 million edges. Graph attention networks aggregate neighborhood information through relation-specific message passing, enabling the model to leverage both direct and indirect associations between genes and phenotypes. This proves critical for rare diseases where only 28\% of patient phenotypic features have known direct associations with causal genes. For patients with phenotypes separated by more than two relationship hops from their causal gene, SHEPHERD achieved 77.8\% accuracy in ranking the correct gene within the top 5 predictions.

Validation on three independent real-world cohorts totaling 2,042 patients with 378 unique genetic conditions demonstrated that synthetic training translates to real-world diagnostic utility. SHEPHERD ranked the correct causal gene first for 40\% of Undiagnosed Diseases Network patients when provided expert-curated candidate gene lists, achieving 2.2-fold improvement over non-guided baselines. Perhaps most remarkably, 81.8\% of real patient phenotype terms appeared in the synthetic training cohort despite the entirely simulated origin, and synthetic patients clustered with real patients of the same disease category in learned embedding spaces.

This knowledge-guided synthesis paradigm differs fundamentally from pure generative modeling. Rather than learning distributions solely from data, the approach encodes structured domain knowledge that constrains generation to biologically plausible combinations. For rare genetic diseases where acquiring thousands of real patients is impossible, this knowledge integration enables deep learning models to achieve clinical utility despite extreme data scarcity. The 20 synthetic patients generated per disease in SHEPHERD's training set would be insufficient for traditional supervised learning but prove adequate when combined with knowledge graph structure and few-shot learning objectives.

Genomic data synthesis extends beyond phenotype-level modeling to include DNA/RNA sequence generation and gene expression profile synthesis. Mendes et al. described applications of GANs and VAEs to capturing rare genetic variants that appear infrequently in population databases \cite{mendes2025synthetic}. CTAB-GAN+ and normalizing flows have generated synthetic patient cohorts for acute myeloid leukemia studies, achieving threefold expansion while preserving the statistical characteristics of rare allelic combinations. These synthetic genomes enable variant calling algorithm development and population genetics analyses without requiring consent protocols or privacy protection mechanisms that constrain real genomic data sharing.

The integration of generative models with knowledge graphs, ontologies, and mechanistic biological models represents a frontier for rare disease applications. Pure data-driven approaches struggle when training examples number in the single or low double digits. Hybrid methods that incorporate known gene-disease relationships, protein-protein interaction networks, metabolic pathway structures, and phenotype ontologies can generate synthetic data that respects biological constraints while expanding available training examples. The balance between data-driven flexibility and knowledge-based constraint remains an active area of investigation, with different applications requiring different trade-offs between these complementary approaches.
